\documentclass[11pt,letterpaper]{article}
\usepackage[T1]{fontenc}
\usepackage{amsmath,amsthm}
\usepackage{newpxtext,newpxmath}
\usepackage[margin=0.95in]{geometry}
\usepackage{microtype,xcolor,xurl}
\usepackage{hyperref}
\definecolor{oliveink}{HTML}{4F6041}
\hypersetup{colorlinks=true,linkcolor=oliveink,urlcolor=oliveink,citecolor=oliveink,
 pdftitle={Genetic gap indicators: the short proof}}
\newcommand{\No}{\mathbf{No}}
\newcommand{\N}{\mathbb{N}}
\newcommand{\R}{\mathbb{R}}
\newcommand{\cut}[2]{\left\{\,#1\;\middle|\;#2\,\right\}}
\newtheorem*{claim}{Claim}
\setlength{\parskip}{0.35em}
\setlength{\emergencystretch}{2em}
\begin{document}
\begin{center}
{\Large\bfseries Genetic gap indicators: the short proof}\par
\vspace{0.2cm}
{\small Research note \;\textbullet\; 20 September 2026}
\end{center}
\noindent\textbf{Targets.} Conjecture~53 of Rubinstein-Salzedo--Swaminathan
\cite{RSS} asserts uniqueness, modulo constants, of genetic primitives.
Chapter~8, Conjecture~3 of Berenbeim \cite{B} asserts the sup property for
entire genetic functions. Here ``genetic'' uses their set-sized option
constructors and the conventional treatment of option expressions that
ignore the options of their input.

\noindent\textbf{Construction.} With $\N=\{0,1,2,\ldots\}$ the ordinary finite
naturals, define
\[
 q(t)=\frac{t}{1+t^2},\qquad
 T(x)=\cut{q(n-x):n\in\N}{},\qquad B(x)=1-T(x).
\]
The denominator is positive, $q(t)$ has the sign of $t$, and
$|q(t)|\le1/2$, since $(|t|-1)^2\ge0$.

\begin{claim}
$B(x)=1$ exactly when $x>n$ for every $n\in\N$, and is $0$ otherwise.
\end{claim}
\begin{proof}
If $x>n$ for all $n$, every left option of $T(x)$ is negative. Thus $0$,
the simplest surreal, lies above all of them, giving $T(x)=0$.
Otherwise choose $n\ge x$. Then one left option is nonnegative, excluding
$0$, while every left option is less than $1$. The simplest admissible
number is now $1$, so $T(x)=1$. Subtract from $1$.
\end{proof}

\noindent\textbf{Geneticity, including the later cofinality condition.}
The globally defined rational $q$ is genetic; the sources explicitly use
the one-option cut $\cut{q(x)}{}$. Each $q(n-x)$ is an admissible affine
translate of a previously constructed function. The option set is fixed
and countable, and the right option set is empty. Thus $T$ is a
nonrecursive genetic adjunction. No input is excluded.

The options reference neither $x^L$ nor $x^R$, so changing the representation
of $x$ leaves the cut unchanged. For every $y<x<z$, the later global
cofinality requirement still reads $q(n-x)<T(x)$. If $T(x)=1$ this follows
from $q(n-x)\le1/2$; if $T(x)=0$ every option is negative. The dual condition
is vacuous. Arithmetic closure makes $B$ and $x-B(x)$ genetic as well.

\noindent\textbf{Derivative.} For every surreal $x$ and every $h$ with
$|h|<1$, membership in the positive infinite class is unchanged by adding
$h$. Indeed, an integer upper bound remains an integer upper bound after
increasing it by $1$; and $x>n+1$ for all $n$ implies $x+h>n$ for all $n$.
Consequently
\[
 B(x+h)=B(x),\qquad
 \frac{B(x+h)-B(x)}h=0\quad(0<|h|<1).
\]
Thus $B'=0$ everywhere in the full surreal $\varepsilon$--$\delta$ sense,
although $B(0)=0$ and $B(\omega)=1$.

\newpage
\noindent\textbf{First counterexample.} Set $F_0(x)=x$ and $F_1(x)=x-B(x)$.
For $0<|h|<1$, the quotient for $F_1$ is exactly $1$, so
\[
 F_0'=F_1'=1,\qquad F_1-F_0=-B\text{ is nonconstant}.
\]
Both functions agree with $x$ at every real argument. The same counterexample
works on the proper bounded open interval $(0,\omega)$, using the points
$1$ and $\omega/2$ to witness nonconstant difference. At each interior
point shrink the derivative neighborhood to stay in the interval.

Moreover, $F_1$ is an increasing bijection. It is the identity below the
finite--infinite gap and translation by $-1$ above it. Translation by a
finite number preserves both sides. Any cross-gap difference is positive
infinite, so subtracting $1$ preserves its sign. The inverse is $x+B(x)$.
The piecewise form of $F_1$ already appears in \cite[p.~37, eq.~(7)]{RSS};
the additional content here is its genetic formula.

\noindent\textbf{Second counterexample.} Consider the sublevel class
\[
 E=\{x\in\No:0<x<\omega,\ B(x)\le1/2\}.
\]
It is exactly the positive finite surreals. Its surreal upper bounds are
exactly the positive infinite surreals: it contains every positive integer,
and each of its members is bounded by an integer. Every such upper bound
$u$ can be replaced by the smaller upper bound $u/2$. Thus $E$ has no
surreal least upper bound. It is nonempty and bounded by $\omega$, so
neither end gap is its supremum. The sup property fails.

\noindent\textbf{Why there is no missing derivative point.} The gap between
numbers bounded above by a natural and numbers above all naturals is not
itself a surreal number. Its lower side has no greatest element and its
upper side has no least element. A derivative is computed at an actual
surreal input, never at this unrealized boundary. In particular, $\omega$
is on the upper side, not at the boundary.

\noindent\textbf{Scope and validation.} These are counterexamples under the
published set-sized option conventions, not a claim about every more
restrictive use of the term ``genetic.'' No later resolution was located
in the recorded searches, but priority is not established. Independent
expert review and proof-assistant verification have not been performed.
The full article contains the general set-indexed construction,
ordinal-indexed independent families, a definition audit, and exact finite
computational checks with their limitations.

\begin{thebibliography}{9}
\bibitem{RSS} S.~Rubinstein-Salzedo and A.~Swaminathan,
\emph{Analysis on Surreal Numbers}, Journal of Logic \& Analysis
\textbf{6}:5 (2014), 1--39. Constructor and examples: pp.~9--11;
Conjecture~53: p.~37.
\url{https://doi.org/10.4115/jla.2014.6.5}.
\bibitem{B} A.~M.~Berenbeim,
\emph{Foundations for the Analysis of Surreal-Valued Genetic Functions},
Ph.D. dissertation, University of Illinois at Chicago, 2022.
Definition~55: pp.~126--128; rational step example: pp.~133--134;
Chapter~8, Conjecture~3: pp.~225--226.
\href{https://indigo.uic.edu/articles/thesis/Foundations_for_the_Analysis_of_Surreal-Valued_Genetic_Functions/20253705/1/files/36188862.pdf}{UIC repository PDF (version 1)}.
\end{thebibliography}
\end{document}
